
\documentclass{article}
\usepackage{colortbl}
\usepackage{makecell}
\usepackage{multirow}
\usepackage{supertabular}

\begin{document}

\newcounter{utterance}

\centering \large Interaction Transcript for game `dond', experiment `coop', episode 0 with gemma{-}4{-}26b{-}a4b{-}it.
\vspace{24pt}

{ \footnotesize  \setcounter{utterance}{1}
\setlength{\tabcolsep}{0pt}
\begin{supertabular}{c@{$\;$}|p{.15\linewidth}@{}p{.15\linewidth}p{.15\linewidth}p{.15\linewidth}p{.15\linewidth}p{.15\linewidth}}
    \# & $\;$A & \multicolumn{4}{c}{Game Master} & $\;\:$B\\
    \hline

    \theutterance \stepcounter{utterance}  
    & & \multicolumn{4}{p{0.6\linewidth}}{
        \cellcolor[rgb]{0.9,0.9,0.9}{
            \makecell[{{p{\linewidth}}}]{
                \texttt{\tiny{[P1$\langle$GM]}}
                \texttt{Вы участвуете в игре кооперативных переговоров, в которой должны договориться с другим игроком о разделении набора предметов между вами.} \\
\\ 
\texttt{Правила:} \\
\texttt{(a) Вы и другой игрок получаете набор предметов. Каждому из вас также дается секретная функция стоимости, представляющая, как вы оцениваете каждый тип объекта.} \\
\texttt{(b) Вы обмениваетесь сообщениями с другим игроком, чтобы согласиться, кто какие предметы получает. Каждый может отправить максимум 5 сообщений или закончить игру раньше, сделав секретное предложение в любой момент.} \\
\texttt{(c) Вас просят отправить секретное предложение, указывающее предметы, которые вы хотите, в следующем формате в квадратных скобках: "{[}Предложение: <число> <названия предмета>, <число> <названия предмета>, <...>{]}"} \\
\texttt{(d) Если ваши предложения дополняют друг друга, то есть есть достаточно предметов для выполнения обоих предложений, каждый игрок получает баллы на основе суммы стоимостей предметов, которые он получил. В противном случае вы оба получаете нулевой балл.} \\
\texttt{(e) Ваша цель {-} максимизировать сумму вашего балла и балла другого игрока.} \\
\\ 
\texttt{Давайте начнем.} \\
\\ 
\texttt{Набор доступных предметов:} \\
\\ 
\texttt{1 болт, 1 сумка, 2 пассатижи, 2 горшки} \\
\\ 
\texttt{Ваша секретная функция стоимости:} \\
\\ 
\texttt{болт: 2, сумка: 2, пассатижи: 0, горшок: 3} \\
\\ 
\texttt{ВАЖНО: Ваши сообщения, если только это не секретное предложение, передаются прямо другому игроку, поэтому не включайте никакие ответы на правила или текст, объявляющий ваше сообщение. Чтобы сделать секретное предложение, используйте указанный формат. Не используйте квадратные скобки при общении с другим игроком, иначе это будет интерпретировано как ваше секретное предложение.} \\
            }
        }
    }
    & \\ \\

    \theutterance \stepcounter{utterance}  
    & \multicolumn{4}{p{0.6\linewidth}}{
        \cellcolor[rgb]{0.9,0.9,0.9}{
            \makecell[{{p{\linewidth}}}]{
                \texttt{\tiny{[P1$\rangle$GM]}}
                \texttt{Привет! Давай договоримся о разделении предметов так, чтобы мы оба получили максимум пользы. Моя ценность предметов: болт (2), сумка (2), пассатижи (0), горшок (3).} \\
\\ 
\texttt{Предлагаю обсудить, какие предметы наиболее важны для тебя, чтобы мы не забирали одно и то же и не получили ноль баллов. Что ты думаешь о таком распределении: я забираю болт, сумку и горшки, а ты забираешь пассатижи? Или у тебя другие приоритеты?} \\
            }
        }
    }
    & & \\ \\

    \theutterance \stepcounter{utterance}  
    & & & \multicolumn{2}{p{0.3\linewidth}}{
        \cellcolor[rgb]{0.9,0.9,0.9}{
            \makecell[{{p{\linewidth}}}]{
                \texttt{\tiny{[GM$|$GM]}}
                \texttt{continue} \\
            }
        }
    }
    & & \\ \\

    \theutterance \stepcounter{utterance}  
    & & & \multicolumn{2}{p{0.3\linewidth}}{
        \cellcolor[rgb]{0.9,0.9,0.9}{
            \makecell[{{p{\linewidth}}}]{
                \texttt{\tiny{[GM$|$GM]}}
                \texttt{Привет! Давай договоримся о разделении предметов так, чтобы мы оба получили максимум пользы. Моя ценность предметов: болт (2), сумка (2), пассатижи (0), горшок (3).} \\
\\ 
\texttt{Предлагаю обсудить, какие предметы наиболее важны для тебя, чтобы мы не забирали одно и то же и не получили ноль баллов. Что ты думаешь о таком распределении: я забираю болт, сумку и горшки, а ты забираешь пассатижи? Или у тебя другие приоритеты?} \\
            }
        }
    }
    & & \\ \\

    \theutterance \stepcounter{utterance}  
    & & \multicolumn{4}{p{0.6\linewidth}}{
        \cellcolor[rgb]{0.9,0.9,0.9}{
            \makecell[{{p{\linewidth}}}]{
                \texttt{\tiny{[GM$\rangle$P2]}}
                \texttt{Вы участвуете в игре кооперативных переговоров, в которой должны договориться с другим игроком о разделении набора предметов между вами.} \\
\\ 
\texttt{Правила:} \\
\texttt{(a) Вы и другой игрок получаете набор предметов. Каждому из вас также дается секретная функция стоимости, представляющая, как вы оцениваете каждый тип объекта.} \\
\texttt{(b) Вы обмениваетесь сообщениями с другим игроком, чтобы согласиться, кто какие предметы получает. Каждый может отправить максимум 5 сообщений или закончить игру раньше, сделав секретное предложение в любой момент.} \\
\texttt{(c) Вас просят отправить секретное предложение, указывающее предметы, которые вы хотите, в следующем формате в квадратных скобках: "{[}Предложение: <число> <названия предмета>, <число> <названия предмета>, <...>{]}"} \\
\texttt{(d) Если ваши предложения дополняют друг друга, то есть есть достаточно предметов для выполнения обоих предложений, каждый игрок получает баллы на основе суммы стоимостей предметов, которые он получил. В противном случае вы оба получаете нулевой балл.} \\
\texttt{(e) Ваша цель {-} максимизировать сумму вашего балла и балла другого игрока.} \\
\\ 
\texttt{Давайте начнем.} \\
\\ 
\texttt{Набор доступных предметов:} \\
\\ 
\texttt{1 болт, 1 сумка, 2 пассатижи, 2 горшки} \\
\\ 
\texttt{Ваша секретная функция стоимости:} \\
\\ 
\texttt{болт: 3, сумка: 1, пассатижи: 3, горшок: 0} \\
\\ 
\texttt{ВАЖНО: Ваши сообщения, если только это не секретное предложение, передаются прямо другому игроку, поэтому не включайте никакие ответы на правила или текст, объявляющий ваше сообщение. Чтобы сделать секретное предложение, используйте указанный формат. Не используйте квадратные скобки при общении с другим игроком, иначе это будет интерпретировано как ваше секретное предложение.} \\
\\ 
\texttt{Привет! Давай договоримся о разделении предметов так, чтобы мы оба получили максимум пользы. Моя ценность предметов: болт (2), сумка (2), пассатижи (0), горшок (3).} \\
\\ 
\texttt{Предлагаю обсудить, какие предметы наиболее важны для тебя, чтобы мы не забирали одно и то же и не получили ноль баллов. Что ты думаешь о таком распределении: я забираю болт, сумку и горшки, а ты забираешь пассатижи? Или у тебя другие приоритеты?} \\
            }
        }
    }
    & \\ \\

    \theutterance \stepcounter{utterance}  
    & & & \multicolumn{4}{p{0.6\linewidth}}{
        \cellcolor[rgb]{0.9,0.9,0.9}{
            \makecell[{{p{\linewidth}}}]{
                \texttt{\tiny{[GM$\langle$P2]}}
                \texttt{Привет! Спасибо за предложение. Давай проанализируем наши ценности, чтобы максимизировать общую сумму.} \\
\\ 
\texttt{Мои приоритеты: болт (3), пассатижи (3), сумка (1), горшок (0).} \\
\texttt{Твои приоритеты: горшок (3), сумка (2), болт (2), пассатижи (0).} \\
\\ 
\texttt{Твое предложение (ты: болт, сумка, горшки; я: пассатижи) дает нам:} \\
\texttt{Твой балл: 2 (болт) + 2 (сумка) + 2*3 (горшки) = 10.} \\
\texttt{Мой балл: 2*0 (пассатижи) = 0.} \\
\texttt{Итого: 10.} \\
\\ 
\texttt{Но если мы разделим предметы по нашим максимальным ценностям, мы можем получить больше. Например:} \\
\texttt{Я забираю болт и пассатижи (3 + 2*3 = 9).} \\
\texttt{Ты забираешь сумку и горшки (2 + 2*3 = 8).} \\
\texttt{Итого: 17.} \\
\\ 
\texttt{Как тебе такой вариант? Я забираю болт и пассатижи, а ты забираешь сумку и горшки. Это кажется более эффективным для общей суммы.} \\
            }
        }
    }
     \\ \\

    \theutterance \stepcounter{utterance}  
    & & & \multicolumn{2}{p{0.3\linewidth}}{
        \cellcolor[rgb]{0.9,0.9,0.9}{
            \makecell[{{p{\linewidth}}}]{
                \texttt{\tiny{[GM$|$GM]}}
                \texttt{continue} \\
            }
        }
    }
    & & \\ \\

    \theutterance \stepcounter{utterance}  
    & & & \multicolumn{2}{p{0.3\linewidth}}{
        \cellcolor[rgb]{0.9,0.9,0.9}{
            \makecell[{{p{\linewidth}}}]{
                \texttt{\tiny{[GM$|$GM]}}
                \texttt{Привет! Спасибо за предложение. Давай проанализируем наши ценности, чтобы максимизировать общую сумму.} \\
\\ 
\texttt{Мои приоритеты: болт (3), пассатижи (3), сумка (1), горшок (0).} \\
\texttt{Твои приоритеты: горшок (3), сумка (2), болт (2), пассатижи (0).} \\
\\ 
\texttt{Твое предложение (ты: болт, сумка, горшки; я: пассатижи) дает нам:} \\
\texttt{Твой балл: 2 (болт) + 2 (сумка) + 2*3 (горшки) = 10.} \\
\texttt{Мой балл: 2*0 (пассатижи) = 0.} \\
\texttt{Итого: 10.} \\
\\ 
\texttt{Но если мы разделим предметы по нашим максимальным ценностям, мы можем получить больше. Например:} \\
\texttt{Я забираю болт и пассатижи (3 + 2*3 = 9).} \\
\texttt{Ты забираешь сумку и горшки (2 + 2*3 = 8).} \\
\texttt{Итого: 17.} \\
\\ 
\texttt{Как тебе такой вариант? Я забираю болт и пассатижи, а ты забираешь сумку и горшки. Это кажется более эффективным для общей суммы.} \\
            }
        }
    }
    & & \\ \\

    \theutterance \stepcounter{utterance}  
    & & \multicolumn{4}{p{0.6\linewidth}}{
        \cellcolor[rgb]{0.9,0.9,0.9}{
            \makecell[{{p{\linewidth}}}]{
                \texttt{\tiny{[P1$\langle$GM]}}
                \texttt{Привет! Спасибо за предложение. Давай проанализируем наши ценности, чтобы максимизировать общую сумму.} \\
\\ 
\texttt{Мои приоритеты: болт (3), пассатижи (3), сумка (1), горшок (0).} \\
\texttt{Твои приоритеты: горшок (3), сумка (2), болт (2), пассатижи (0).} \\
\\ 
\texttt{Твое предложение (ты: болт, сумка, горшки; я: пассатижи) дает нам:} \\
\texttt{Твой балл: 2 (болт) + 2 (сумка) + 2*3 (горшки) = 10.} \\
\texttt{Мой балл: 2*0 (пассатижи) = 0.} \\
\texttt{Итого: 10.} \\
\\ 
\texttt{Но если мы разделим предметы по нашим максимальным ценностям, мы можем получить больше. Например:} \\
\texttt{Я забираю болт и пассатижи (3 + 2*3 = 9).} \\
\texttt{Ты забираешь сумку и горшки (2 + 2*3 = 8).} \\
\texttt{Итого: 17.} \\
\\ 
\texttt{Как тебе такой вариант? Я забираю болт и пассатижи, а ты забираешь сумку и горшки. Это кажется более эффективным для общей суммы.} \\
            }
        }
    }
    & \\ \\

    \theutterance \stepcounter{utterance}  
    & \multicolumn{4}{p{0.6\linewidth}}{
        \cellcolor[rgb]{0.9,0.9,0.9}{
            \makecell[{{p{\linewidth}}}]{
                \texttt{\tiny{[P1$\rangle$GM]}}
                \texttt{Твой расчет очень точный, и я полностью согласен! Твой вариант действительно значительно увеличивает общую сумму баллов (17 против 10).} \\
\\ 
\texttt{Давай проверим еще раз:} \\
\texttt{Если ты берешь болт (3) и 2 пассатижи (2 * 3 = 6), твой балл будет 9.} \\
\texttt{Если я беру сумку (2) и 2 горшка (2 * 3 = 6), мой балл будет 8.} \\
\texttt{Итого: 17.} \\
\\ 
\texttt{Все предметы распределены, и наши предложения дополняют друг друга. Я согласен на этот вариант. Если ты не против, я готов отправить секретное предложение.} \\
            }
        }
    }
    & & \\ \\

    \theutterance \stepcounter{utterance}  
    & & & \multicolumn{2}{p{0.3\linewidth}}{
        \cellcolor[rgb]{0.9,0.9,0.9}{
            \makecell[{{p{\linewidth}}}]{
                \texttt{\tiny{[GM$|$GM]}}
                \texttt{continue} \\
            }
        }
    }
    & & \\ \\

    \theutterance \stepcounter{utterance}  
    & & & \multicolumn{2}{p{0.3\linewidth}}{
        \cellcolor[rgb]{0.9,0.9,0.9}{
            \makecell[{{p{\linewidth}}}]{
                \texttt{\tiny{[GM$|$GM]}}
                \texttt{Твой расчет очень точный, и я полностью согласен! Твой вариант действительно значительно увеличивает общую сумму баллов (17 против 10).} \\
\\ 
\texttt{Давай проверим еще раз:} \\
\texttt{Если ты берешь болт (3) и 2 пассатижи (2 * 3 = 6), твой балл будет 9.} \\
\texttt{Если я беру сумку (2) и 2 горшка (2 * 3 = 6), мой балл будет 8.} \\
\texttt{Итого: 17.} \\
\\ 
\texttt{Все предметы распределены, и наши предложения дополняют друг друга. Я согласен на этот вариант. Если ты не против, я готов отправить секретное предложение.} \\
            }
        }
    }
    & & \\ \\

    \theutterance \stepcounter{utterance}  
    & & \multicolumn{4}{p{0.6\linewidth}}{
        \cellcolor[rgb]{0.9,0.9,0.9}{
            \makecell[{{p{\linewidth}}}]{
                \texttt{\tiny{[GM$\rangle$P2]}}
                \texttt{Твой расчет очень точный, и я полностью согласен! Твой вариант действительно значительно увеличивает общую сумму баллов (17 против 10).} \\
\\ 
\texttt{Давай проверим еще раз:} \\
\texttt{Если ты берешь болт (3) и 2 пассатижи (2 * 3 = 6), твой балл будет 9.} \\
\texttt{Если я беру сумку (2) и 2 горшка (2 * 3 = 6), мой балл будет 8.} \\
\texttt{Итого: 17.} \\
\\ 
\texttt{Все предметы распределены, и наши предложения дополняют друг друга. Я согласен на этот вариант. Если ты не против, я готов отправить секретное предложение.} \\
            }
        }
    }
    & \\ \\

    \theutterance \stepcounter{utterance}  
    & & & \multicolumn{4}{p{0.6\linewidth}}{
        \cellcolor[rgb]{0.9,0.9,0.9}{
            \makecell[{{p{\linewidth}}}]{
                \texttt{\tiny{[GM$\langle$P2]}}
                \texttt{Я полностью согласен! Это идеальное распределение, которое максимизирует нашу общую выгоду. Давай закрепим это.} \\
\\ 
\texttt{{[}Предложение: 1 болт, 2 пассатижи{]}} \\
            }
        }
    }
     \\ \\

    \theutterance \stepcounter{utterance}  
    & & & \multicolumn{2}{p{0.3\linewidth}}{
        \cellcolor[rgb]{0.9,0.9,0.9}{
            \makecell[{{p{\linewidth}}}]{
                \texttt{\tiny{[GM$|$GM]}}
                \texttt{proposal} \\
            }
        }
    }
    & & \\ \\

    \theutterance \stepcounter{utterance}  
    & & & \multicolumn{2}{p{0.3\linewidth}}{
        \cellcolor[rgb]{0.9,0.9,0.9}{
            \makecell[{{p{\linewidth}}}]{
                \texttt{\tiny{[GM$|$GM]}}
                [1, 0, 2, 0]
            }
        }
    }
    & & \\ \\

    \theutterance \stepcounter{utterance}  
    & & \multicolumn{4}{p{0.6\linewidth}}{
        \cellcolor[rgb]{0.9,0.9,0.9}{
            \makecell[{{p{\linewidth}}}]{
                \texttt{\tiny{[P1$\langle$GM]}}
                \texttt{Другой игрок отправил секретное предложение. Теперь ваша очередь отправить своё собственное предложение.} \\
            }
        }
    }
    & \\ \\

    \theutterance \stepcounter{utterance}  
    & \multicolumn{4}{p{0.6\linewidth}}{
        \cellcolor[rgb]{0.9,0.9,0.9}{
            \makecell[{{p{\linewidth}}}]{
                \texttt{\tiny{[P1$\rangle$GM]}}
                \texttt{{[}Предложение: 1 сумка, 2 горшка{]}} \\
            }
        }
    }
    & & \\ \\

    \theutterance \stepcounter{utterance}  
    & & & \multicolumn{2}{p{0.3\linewidth}}{
        \cellcolor[rgb]{0.9,0.9,0.9}{
            \makecell[{{p{\linewidth}}}]{
                \texttt{\tiny{[GM$|$GM]}}
                \texttt{proposal} \\
            }
        }
    }
    & & \\ \\

    \theutterance \stepcounter{utterance}  
    & & & \multicolumn{2}{p{0.3\linewidth}}{
        \cellcolor[rgb]{0.9,0.9,0.9}{
            \makecell[{{p{\linewidth}}}]{
                \texttt{\tiny{[GM$|$GM]}}
                [0, 1, 0, 2]
            }
        }
    }
    & & \\ \\

    \theutterance \stepcounter{utterance}  
    & & & \multicolumn{2}{p{0.3\linewidth}}{
        \cellcolor[rgb]{0.9,0.9,0.9}{
            \makecell[{{p{\linewidth}}}]{
                \texttt{\tiny{[GM$|$GM]}}
                [[0, 1, 0, 2], [1, 0, 2, 0]]
            }
        }
    }
    & & \\ \\

\end{supertabular}
}

\end{document}
